\documentclass[letterpaper,12pt]{article}

\usepackage{fancyhdr}
\usepackage[most]{tcolorbox}
\usepackage{longtable}
\usepackage{makecell}
\usepackage{tikz}
\usepackage{geometry}
\usepackage{array}
\usepackage{multicol}
\usepackage[table]{xcolor}
\usepackage{polyglossia}
\setmainlanguage{english}
\setotherlanguage{arabic}
\setotherlanguage{hebrew}
\usepackage{fontspec}
\setmainfont{Charis SIL}
\newfontfamily\arabicfont[Script=Arabic,Scale=1.2]{Amiri}
\newfontfamily\hebrewfont{Ezra SIL}[Script=Hebrew]
\geometry{margin=1.5cm}
\usepackage{booktabs}
\usepackage{graphicx}
\usepackage{multirow}
\usepackage{titlesec}
\titlespacing*{\section}
  {0pt}                    % left margin
  {3.5ex plus 2ex minus 1.5ex}  % space before (more flexible)
  {2.3ex plus 1ex minus 0.5ex}  % space after (more flexible)

\setlength{\headheight}{14.49998pt}
\addtolength{\topmargin}{-2.49998pt}
\pagestyle{fancy}
\fancyhf{}
\fancyfoot[C]{\thepage}
\fancyhead[R]{\textit{Are We in the Abbasid or Umayyad Era?}}
\renewcommand{\headrulewidth}{0pt}

\usetikzlibrary{shapes,arrows,decorations.pathmorphing}

% Custom colors
\definecolor{headercolor}{RGB}{70,130,180}
\definecolor{boxcolor}{RGB}{240,248,255}
\definecolor{accentcolor}{RGB}{220,20,60}
\definecolor{tableheader}{RGB}{220,220,220}
\definecolor{dialectcolor}{RGB}{34,139,34}

\begin{document}

\title{\textbf{\Large Arabic Phrase Analyser}\\
\large Are We in the Abbasid Era or the Umayyad State Era?\\
\normalsize \textarabic{هل نحن بعصر عباسي أم عصر الدولة الأموية}}
\author{Igor Deruga}
\date{}
\maketitle

% ======================== Phrase Display ========================
\begin{tcolorbox}[colback=boxcolor,colframe=headercolor,title=\textbf{Arabic Phrase with Full Diacritics},breakable]
\centering
\Large\textarabic{هَلْ نَحْنُ بِعَصْرٍ عَبَّاسِيٍّ أَمْ عَصْرِ الدَّوْلَةِ الْأُمَوِيَّةِ؟}
\end{tcolorbox}

% ======================== Translations ========================
\section{English Translation}
\begin{tcolorbox}[colback=white,colframe=accentcolor,breakable]
\textbf{Literal:} Are we in-era Abbasid or era the-state the-Umayyad?\\
\textit{[Arabic order retained for direct mapping]}\\[0.5em]
\textbf{Adapted:} Are we in an Abbasid era or the era of the Umayyad state?
\end{tcolorbox}

% ======================== Detailed Word Analysis ========================
\section{Detailed Word Analysis}

\subsection{\textarabic{هَلْ} — /hal/}
\begin{tabular}{p{3cm}p{10cm}}
\toprule
\textbf{Translation} & (question particle indicating yes/no question) \\
\textbf{Root} & Interrogative particle (not from triliteral root) \\
\textbf{Pattern} & Independent interrogative particle \\
\textbf{Grammar} & Interrogative particle introducing polar (yes/no) questions \\
\midrule
\textbf{Examples} & \makecell[l]{\parbox{9.5cm}{
1. \textarabic{هَلْ أَنْتَ مُسْتَعِدٌّ؟} — Are you ready? /hal ʔanta mustaʕiddun/\\
2. \textarabic{هَلْ فَهِمْتَ الدَّرْسَ؟} — Did you understand the lesson? /hal fahimta ddarsa/\\
3. \textarabic{هَلْ يُمْكِنُكَ الْمَجِيءُ؟} — Can you come? /hal jumkinuka lmadʒiːʔu/
}} \\
\midrule
\textbf{Function} & Requires yes (\textarabic{نَعَمْ}) or no (\textarabic{لَا}) answer \\
\textbf{Etymology} & From Proto-Semitic interrogative particle; Hebrew \texthebrew{הַ} (ha-) question prefix \\
\bottomrule
\end{tabular}

\subsubsection*{Usage Notes}
\begin{itemize}
  \item \textarabic{هَلْ} introduces binary questions expecting yes/no answers
  \item Different from \textarabic{مَا} (what) and \textarabic{أَ} (general interrogative particle)
  \item Can be used with nominal or verbal sentences
  \item Often indicates genuine inquiry or rhetorical questioning
\end{itemize}

\subsection{\textarabic{نَحْنُ} — /naħnu/}
\begin{tabular}{p{3cm}p{10cm}}
\toprule
\textbf{Translation} & we \\
\textbf{Root} & Independent pronoun (not from triliteral root) \\
\textbf{Pattern} & First person plural independent pronoun \\
\textbf{Grammar} & Independent pronoun, nominative case (subject of nominal sentence) \\
\midrule
\textbf{Examples} & \makecell[l]{\parbox{9.5cm}{
1. \textarabic{نَحْنُ طُلَّابٌ} — We are students /naħnu tˤullaːbun/\\
2. \textarabic{نَحْنُ نَدْرُسُ الْعَرَبِيَّةَ} — We study Arabic /naħnu nadrusu lʕarabijjata/\\
3. \textarabic{نَحْنُ هُنَا} — We are here /naħnu hunaː/
}} \\
\midrule
\textbf{Related Forms} & \textarabic{نَا} (attached suffix: our/us), \textarabic{إِيَّانَا} (emphatic: us) \\
\textbf{Etymology} & From Proto-Semitic *naħnu; Hebrew \texthebrew{אֲנַחְנוּ} (ʔanaħnu) "we" \\
\bottomrule
\end{tabular}

\subsubsection*{Pronoun System}
\begin{tabular}{|c|c|c|}
\hline
\textbf{Person} & \textbf{Singular} & \textbf{Plural} \\
\hline
1st person & \textarabic{أَنَا} /ʔanaː/ (I) & \textarabic{نَحْنُ} /naħnu/ (we) \\
\hline
2nd person masc. & \textarabic{أَنْتَ} /ʔanta/ (you) & \textarabic{أَنْتُمْ} /ʔantum/ (you) \\
\hline
2nd person fem. & \textarabic{أَنْتِ} /ʔanti/ (you) & \textarabic{أَنْتُنَّ} /ʔantunna/ (you) \\
\hline
3rd person masc. & \textarabic{هُوَ} /huwa/ (he) & \textarabic{هُمْ} /hum/ (they) \\
\hline
3rd person fem. & \textarabic{هِيَ} /hija/ (she) & \textarabic{هُنَّ} /hunna/ (they) \\
\hline
\end{tabular}

\subsection{\textarabic{بِعَصْرٍ} — /biʕasˤrin/}
\begin{tabular}{p{3cm}p{10cm}}
\toprule
\textbf{Translation} & in an era / in a period \\
\textbf{Root} & \textarabic{ع-ص-ر} (ʕ-sˤ-r) \\
\textbf{Pattern} & \textarabic{فَعْل} (faʕl) with preposition \textarabic{بِ} \\
\textbf{Grammar} & Preposition \textarabic{بِ} + masculine noun, genitive case (after preposition), indefinite \\
\midrule
\textbf{Examples} & \makecell[l]{\parbox{9.5cm}{
1. \textarabic{عَصْرُ النَّهْضَةِ} — The Renaissance era /ʕasˤru nnahðˤati/\\
2. \textarabic{فِي عَصْرِنَا} — In our era /fiː ʕasˤrinaː/\\
3. \textarabic{عَصْرٌ جَدِيدٌ} — A new era /ʕasˤrun dʒadiːdun/
}} \\
\midrule
\textbf{Synonyms} & \textarabic{زَمَن} (time), \textarabic{دَهْر} (age), \textarabic{حِقْبَة} (period), \textarabic{عَهْد} (epoch) \\
\textbf{Etymology} & From root meaning "to press, squeeze"; metaphorically extended to "time period" \\
\bottomrule
\end{tabular}

\subsubsection*{Declension}
\begin{tabular}{|c|c|c|c|}
\hline
\textbf{Case} & \textbf{Indefinite} & \textbf{Definite} & \textbf{With \textarabic{بِ}} \\
\hline
Nominative & \textarabic{عَصْرٌ} /ʕasˤrun/ & \textarabic{الْعَصْرُ} /alʕasˤru/ & — \\
\hline
Accusative & \textarabic{عَصْرًا} /ʕasˤran/ & \textarabic{الْعَصْرَ} /alʕasˤra/ & — \\
\hline
Genitive & \textarabic{عَصْرٍ} /ʕasˤrin/ & \textarabic{الْعَصْرِ} /alʕasˤri/ & \textarabic{بِعَصْرٍ} /biʕasˤrin/ \\
\hline
\end{tabular}

\subsubsection*{Preposition \textarabic{بِ}}
The preposition \textarabic{بِ} /bi/ has multiple meanings:
\begin{itemize}
\item \textbf{Location:} in, at — \textarabic{بِالْبَيْتِ} (at home)
\item \textbf{Instrument:} with, by means of — \textarabic{كَتَبَ بِالْقَلَمِ} (wrote with the pen)
\item \textbf{Association:} with — \textarabic{مَرْحَبًا بِكَ} (welcome to you)
\item \textbf{Temporal:} in (time) — \textarabic{بِعَصْرٍ} (in an era)
\end{itemize}

\subsection{\textarabic{عَبَّاسِيٍّ} — /ʕabbaːsijjin/}
\begin{tabular}{p{3cm}p{10cm}}
\toprule
\textbf{Translation} & Abbasid (relating to the Abbasid dynasty/caliphate) \\
\textbf{Root} & \textarabic{ع-ب-س} (ʕ-b-s) from proper name \textarabic{عَبَّاس} \\
\textbf{Pattern} & \textarabic{فَعَّالِيّ} (faʕʕaːlijj) — nisba (relational) adjective \\
\textbf{Grammar} & Nisba adjective, masculine singular, genitive case (agreeing with \textarabic{عَصْرٍ}), indefinite \\
\midrule
\textbf{Examples} & \makecell[l]{\parbox{9.5cm}{
1. \textarabic{الْخِلَافَةُ الْعَبَّاسِيَّةُ} — The Abbasid Caliphate /alxilaːfatu lʕabbaːsijjatu/\\
2. \textarabic{الْعَصْرُ الْعَبَّاسِيُّ الْأَوَّلُ} — The first Abbasid era /alʕasˤru lʕabbaːsijju l-ʔawwalu/\\
3. \textarabic{شَاعِرٌ عَبَّاسِيٌّ} — An Abbasid poet /ʃaːʕirun ʕabbaːsijjun/
}} \\
\midrule
\textbf{Historical Context} & The Abbasid Caliphate (750-1258 CE), known for cultural flourishing, Baghdad as capital, Golden Age of Islam \\
\textbf{Etymology} & From \textarabic{عَبَّاس} (ʕAbbaːs), uncle of Prophet Muhammad; dynasty claimed descent from him \\
\bottomrule
\end{tabular}

\subsubsection*{Nisba Adjective Formation}
Nisba adjectives relate nouns to places, tribes, or concepts by adding \textarabic{ـِيّ} /-ijj/:

\begin{tabular}{|l|l|l|}
\hline
\textbf{Base Noun} & \textbf{Nisba Adjective} & \textbf{Meaning} \\
\hline
\textarabic{عَبَّاس} /ʕabbaːs/ & \textarabic{عَبَّاسِيّ} /ʕabbaːsijj/ & Abbasid \\
\hline
\textarabic{أُمَيَّة} /ʔumajja/ & \textarabic{أُمَوِيّ} /ʔumawijj/ & Umayyad \\
\hline
\textarabic{مِصْر} /misˤr/ & \textarabic{مِصْرِيّ} /misˤrijj/ & Egyptian \\
\hline
\textarabic{عَرَب} /ʕarab/ & \textarabic{عَرَبِيّ} /ʕarabijj/ & Arab/Arabic \\
\hline
\end{tabular}

\subsubsection*{Declension of Nisba Adjectives}
\begin{tabular}{|c|c|c|}
\hline
\textbf{Case} & \textbf{Indefinite} & \textbf{Definite} \\
\hline
Nominative & \textarabic{عَبَّاسِيٌّ} /ʕabbaːsijjun/ & \textarabic{الْعَبَّاسِيُّ} /alʕabbaːsijju/ \\
\hline
Accusative & \textarabic{عَبَّاسِيًّا} /ʕabbaːsijjan/ & \textarabic{الْعَبَّاسِيَّ} /alʕabbaːsijja/ \\
\hline
Genitive & \textarabic{عَبَّاسِيٍّ} /ʕabbaːsijjin/ & \textarabic{الْعَبَّاسِيِّ} /alʕabbaːsijji/ \\
\hline
\end{tabular}

\subsection{\textarabic{أَمْ} — /ʔam/}
\begin{tabular}{p{3cm}p{10cm}}
\toprule
\textbf{Translation} & or (in questions) \\
\textbf{Root} & Disjunctive particle (not from triliteral root) \\
\textbf{Pattern} & Interrogative disjunction particle \\
\textbf{Grammar} & Coordinating conjunction used in alternative questions \\
\midrule
\textbf{Examples} & \makecell[l]{\parbox{9.5cm}{
1. \textarabic{أَأَنْتَ مُحَمَّدٌ أَمْ عَلِيٌّ؟} — Are you Muhammad or Ali? /ʔaʔanta muħammadun ʔam ʕalijjun/\\
2. \textarabic{هَلْ تُرِيدُ شَايًا أَمْ قَهْوَةً؟} — Do you want tea or coffee? /hal turiːdu ʃaːjan ʔam qahwatan/\\
3. \textarabic{أَنَذْهَبُ أَمْ نَبْقَى؟} — Do we go or stay? /ʔanaðhabu ʔam nabqaː/
}} \\
\midrule
\textbf{Comparison} & Different from \textarabic{أَوْ} (or) which is used in statements, not questions \\
\textbf{Etymology} & From Proto-Semitic interrogative particle \\
\bottomrule
\end{tabular}

\subsubsection*{Usage: \textarabic{أَمْ} vs. \textarabic{أَوْ}}
\begin{tabular}{|l|p{5cm}|p{5cm}|}
\hline
\textbf{Particle} & \textbf{Usage} & \textbf{Example} \\
\hline
\textarabic{أَمْ} /ʔam/ & In questions (alternative) & \textarabic{هَلْ هَذَا أَمْ ذَاكَ؟} (Is it this or that?) \\
\hline
\textarabic{أَوْ} /ʔaw/ & In statements (choice) & \textarabic{خُذْ هَذَا أَوْ ذَاكَ} (Take this or that) \\
\hline
\end{tabular}

\subsection{\textarabic{عَصْرِ} — /ʕasˤri/}
\begin{tabular}{p{3cm}p{10cm}}
\toprule
\textbf{Translation} & era (of) / period (of) \\
\textbf{Root} & \textarabic{ع-ص-ر} (ʕ-sˤ-r) \\
\textbf{Pattern} & \textarabic{فَعْل} (faʕl) \\
\textbf{Grammar} & Masculine noun, genitive case (first term of iḍāfah), definite through iḍāfah \\
\midrule
\textbf{Function} & First term (muḍāf) in construct state with \textarabic{الدَّوْلَةِ} \\
\textbf{Note} & Same word as \textarabic{بِعَصْرٍ} earlier, but here in iḍāfah construction \\
\bottomrule
\end{tabular}

\subsubsection*{Iḍāfah Construction}
\textarabic{عَصْرِ الدَّوْلَةِ الْأُمَوِيَّةِ} /ʕasˤri ddawlati l-ʔumawijjati/

\begin{tabular}{|l|l|l|}
\hline
\textbf{Component} & \textbf{Role} & \textbf{Case} \\
\hline
\textarabic{عَصْرِ} & Muḍāf (possessed) & Genitive (from syntax) \\
\hline
\textarabic{الدَّوْلَةِ} & Muḍāf ilayh & Genitive (required) \\
\hline
\textarabic{الْأُمَوِيَّةِ} & Adjective & Genitive (agreement) \\
\hline
\end{tabular}

\subsection{\textarabic{الدَّوْلَةِ} — /ddawlati/}
\begin{tabular}{p{3cm}p{10cm}}
\toprule
\textbf{Translation} & the state / the dynasty / the empire \\
\textbf{Root} & \textarabic{د-و-ل} (d-w-l) \\
\textbf{Pattern} & \textarabic{فَعْلَة} (faʕlah) \\
\textbf{Grammar} & Feminine noun, genitive case (second term of iḍāfah), definite \\
\midrule
\textbf{Examples} & \makecell[l]{\parbox{9.5cm}{
1. \textarabic{الدَّوْلَةُ الْإِسْلَامِيَّةُ} — The Islamic state /addawlatu l-ʔislaːmijjatu/\\
2. \textarabic{رَئِيسُ الدَّوْلَةِ} — Head of state /raʔiːsu ddawlati/\\
3. \textarabic{دَوْلَةٌ قَوِيَّةٌ} — A strong state /dawlatun qawijjatun/
}} \\
\midrule
\textbf{Synonyms} & \textarabic{حُكُومَة} (government), \textarabic{سُلْطَة} (authority), \textarabic{مَمْلَكَة} (kingdom) \\
\textbf{Etymology} & From root meaning "to alternate, rotate"; extended to mean "ruling power, dynasty" \\
\bottomrule
\end{tabular}

\subsubsection*{Declension}
\begin{tabular}{|c|c|c|}
\hline
\textbf{Case} & \textbf{Indefinite} & \textbf{Definite} \\
\hline
Nominative & \textarabic{دَوْلَةٌ} /dawlatun/ & \textarabic{الدَّوْلَةُ} /addawlatu/ \\
\hline
Accusative & \textarabic{دَوْلَةً} /dawlatan/ & \textarabic{الدَّوْلَةَ} /addawlata/ \\
\hline
Genitive & \textarabic{دَوْلَةٍ} /dawlatin/ & \textarabic{الدَّوْلَةِ} /addawlati/ \\
\hline
\end{tabular}

\subsubsection*{Related Verbal Forms}
\par{\large \textbf{Verbal Noun}: \textarabic{دَوْلَة} /dawlah/ or \textarabic{تَدَاوُل} /tadaːwul/}

\begin{longtable}{|>{\raggedright}p{3.5cm}|p{5cm}|p{5cm}|}
\hline
\textbf{Person} & \textbf{Perfect (Past)} & \textbf{Imperfect (Present)} \\
\hline
\textbf{3rd person masc. sing.} & \textarabic{دَالَ} /daːla/ & \textarabic{يَدُولُ} /jaduːlu/ \\
\hline
\textbf{3rd person fem. sing.} & \textarabic{دَالَتْ} /daːlat/ & \textarabic{تَدُولُ} /taduːlu/ \\
\hline
\textbf{3rd person masc. plural} & \textarabic{دَالُوا} /daːluː/ & \textarabic{يَدُولُونَ} /jaduːluːna/ \\
\hline
\textbf{2nd person masc. sing.} & \textarabic{دُلْتَ} /dulta/ & \textarabic{تَدُولُ} /taduːlu/ \\
\hline
\textbf{1st person singular} & \textarabic{دُلْتُ} /dultu/ & \textarabic{أَدُولُ} /ʔaduːlu/ \\
\hline
\end{longtable}

\subsection{\textarabic{الْأُمَوِيَّةِ} — /l-ʔumawijjati/}
\begin{tabular}{p{3cm}p{10cm}}
\toprule
\textbf{Translation} & the Umayyad (relating to the Umayyad dynasty) \\
\textbf{Root} & \textarabic{أ-م-ي} from proper name \textarabic{أُمَيَّة} \\
\textbf{Pattern} & \textarabic{فُعَلِيَّة} (fuʕalijjah) — feminine nisba adjective \\
\textbf{Grammar} & Nisba adjective, feminine singular, genitive case (agreeing with \textarabic{الدَّوْلَةِ}), definite \\
\midrule
\textbf{Examples} & \makecell[l]{\parbox{9.5cm}{
1. \textarabic{الْخِلَافَةُ الْأُمَوِيَّةُ} — The Umayyad Caliphate /alxilaːfatu l-ʔumawijjatu/\\
2. \textarabic{الْمَسْجِدُ الْأُمَوِيُّ} — The Umayyad Mosque /almasdʒidu l-ʔumawijju/\\
3. \textarabic{عَصْرٌ أُمَوِيٌّ} — An Umayyad era /ʕasˤrun ʔumawijjun/
}} \\
\midrule
\textbf{Historical Context} & The Umayyad Caliphate (661-750 CE), first major Islamic dynasty, capital in Damascus \\
\textbf{Etymology} & From \textarabic{أُمَيَّة} (Umayya), ancestor of the Umayyad clan of Quraysh tribe \\
\bottomrule
\end{tabular}

\subsubsection*{Declension of Feminine Nisba}
\begin{tabular}{|c|c|c|}
\hline
\textbf{Case} & \textbf{Indefinite} & \textbf{Definite} \\
\hline
Nominative & \textarabic{أُمَوِيَّةٌ} /ʔumawijjatun/ & \textarabic{الْأُمَوِيَّةُ} /al-ʔumawijjatu/ \\
\hline
Accusative & \textarabic{أُمَوِيَّةً} /ʔumawijjatan/ & \textarabic{الْأُمَوِيَّةَ} /al-ʔumawijjata/ \\
\hline
Genitive & \textarabic{أُمَوِيَّةٍ} /ʔumawijjatin/ & \textarabic{الْأُمَوِيَّةِ} /al-ʔumawijjati/ \\
\hline
\end{tabular}

\subsubsection*{Historical Comparison}
\begin{tabular}{|l|p{5cm}|p{5cm}|}
\hline
\textbf{Feature} & \textbf{Umayyad (661-750)} & \textbf{Abbasid (750-1258)} \\
\hline
Capital & Damascus & Baghdad \\
\hline
Founder & Muʕāwiyah ibn Abī Sufyān & Abū al-ʕAbbās as-Saffāḥ \\
\hline
Character & Arab aristocracy & Multi-ethnic, Persian influence \\
\hline
Legacy & Expansion, Arabic spread & Golden Age, sciences, culture \\
\hline
\end{tabular}

\section{Phrase Analysis}
\noindent
\begin{tcolorbox}[colback=boxcolor,colframe=headercolor,breakable]
\textbf{Grammatical Structure:}\\
Interrogative particle + independent pronoun (subject) + prepositional phrase with indefinite noun + nisba adjective (attribute) + disjunctive particle + iḍāfah construction (definite noun + definite noun + definite nisba adjective)

\vspace{0.5em}
\textbf{Sentence Type:} Nominal sentence (\textarabic{جُمْلَة اسْمِيَّة}) introduced by yes/no interrogative \textarabic{هَلْ}

\vspace{0.5em}
\textbf{Key Grammar Points:}
\begin{itemize}
\item \textbf{Yes/no question:} \textarabic{هَلْ} introduces alternative question expecting choice between two options
\item \textbf{Subject:} \textarabic{نَحْنُ} (we) — independent pronoun functioning as subject
\item \textbf{First predicate:} \textarabic{بِعَصْرٍ عَبَّاسِيٍّ} — prepositional phrase with indefinite noun modified by nisba adjective
\item \textbf{Disjunction:} \textarabic{أَمْ}— "or" introduces second alternative
\item \textbf{Second predicate:} \textarabic{عَصْرِ الدَّوْلَةِ الْأُمَوِيَّةِ} — iḍāfah construction (definite through association)
\item \textbf{Asymmetry:} First option indefinite (\textarabic{بِعَصْرٍ عَبَّاسِيٍّ}), second definite (\textarabic{عَصْرِ الدَّوْلَةِ الْأُمَوِيَّةِ})
\item \textbf{Ellipsis:} The preposition \textarabic{بِ} is understood before the second predicate (full form: \textarabic{أَمْ بِعَصْرِ الدَّوْلَةِ...})
\item \textbf{Nisba agreement:} Both \textarabic{عَبَّاسِيٍّ} and \textarabic{الْأُمَوِيَّةِ} agree with their respective nouns in all features
\item \textbf{Rhetorical function:} Question implies anachronistic situation or criticism of outdated governance
\item \textbf{Historical reference:} Contrasts two major Islamic dynasties (Umayyad 661-750 CE vs. Abbasid 750-1258 CE)
\end{itemize}
\end{tcolorbox}

\section{Syntactic Analysis}

\begin{tcolorbox}[colback=white,colframe=accentcolor,breakable]
\textbf{Phrase Structure Tree:}

\begin{center}
\begin{tikzpicture}[
  level distance=1.2cm,
  level 1/.style={sibling distance=7cm},
  level 2/.style={sibling distance=3.5cm},
  level 3/.style={sibling distance=2.5cm},
  level 4/.style={sibling distance=2cm},
  every node/.style={rectangle, draw, rounded corners, minimum width=2.2cm, align=center, font=\small}
]

\node {Yes/No Question}
  child {node {\textarabic{هَلْ}\\(Are?)}}
  child {node {Nominal Sentence}
    child {node {Subject\\\textarabic{نَحْنُ}\\(we)}}
    child {node {Alternative\\Predicates}
      child {node {Predicate 1}
        child {node {\textarabic{بِعَصْرٍ}\\in era}}
        child {node {\textarabic{عَبَّاسِيٍّ}\\Abbasid}}
      }
      child {node {\textarabic{أَمْ}\\(or)}}
      child {node {Predicate 2}
        child {node {Iḍāfah\\\textarabic{عَصْرِ الدَّوْلَةِ}\\era of state}}
        child {node {\textarabic{الْأُمَوِيَّةِ}\\the Umayyad}}
      }
    }
  };

\end{tikzpicture}
\end{center}

\vspace{1em}
\textbf{Syntactic Relationships:}
\begin{enumerate}
\item \textbf{Main clause:} \textarabic{هَلْ نَحْنُ...} — Yes/no interrogative nominal sentence
\item \textbf{Subject-predicate:} \textarabic{نَحْنُ} (subject) + two alternative predicates joined by \textarabic{أَمْ}
\item \textbf{Predicate 1 structure:} Prepositional phrase \textarabic{بِعَصْرٍ} + attributive adjective \textarabic{عَبَّاسِيٍّ}
\item \textbf{Predicate 2 structure:} Iḍāfah chain \textarabic{عَصْرِ الدَّوْلَةِ} + attributive adjective \textarabic{الْأُمَوِيَّةِ}
\item \textbf{Ellipsis:} Preposition \textarabic{بِ} omitted before second predicate (common in alternative constructions)
\item \textbf{Definiteness contrast:} First predicate indefinite (any Abbasid era), second definite (the specific Umayyad state era)
\end{enumerate}

\vspace{1em}
\textbf{Simplified Box Diagram:}

\begin{center}
\begin{tikzpicture}[
  box/.style={rectangle, draw, minimum width=2.5cm, minimum height=0.8cm, align=center},
  arrow/.style={->, >=stealth, thick}
]

% Top level
\node[box] (hal) at (0,0) {\textarabic{هَلْ}\\Are?};
\node[box] (nahnu) at (3,0) {\textarabic{نَحْنُ}\\we};

% First alternative
\node[box] (bi3asr) at (1,-2) {\textarabic{بِعَصْرٍ}\\in an era};
\node[box] (3abbasi) at (4,-2) {\textarabic{عَبَّاسِيٍّ}\\Abbasid};

% Disjunction
\node[box] (am) at (7,-2) {\textarabic{أَمْ}\\or};

% Second alternative
\node[box] (3asr) at (7,-4) {\textarabic{عَصْرِ}\\era (of)};
\node[box] (dawla) at (10,-4) {\textarabic{الدَّوْلَةِ}\\the state};
\node[box] (umawi) at (10,-6) {\textarabic{الْأُمَوِيَّةِ}\\the Umayyad};

% Arrows
\draw[arrow] (hal) -- (nahnu);
\draw[arrow] (nahnu) -- (bi3asr);
\draw[arrow] (bi3asr) -- (3abbasi);
\draw[arrow] (nahnu) -- (am);
\draw[arrow] (am) -- (3asr);
\draw[arrow] (3asr) -- (dawla);
\draw[arrow] (dawla) -- (umawi);

\end{tikzpicture}
\end{center}
\end{tcolorbox}

% ======================== Similar Phrases for Practice ========================
\section{Similar Phrases for Practice}

\begin{enumerate}
\item \textarabic{هَلْ نَحْنُ فِي زَمَنٍ إِسْلَامِيٍّ أَمْ زَمَنِ الْجَاهِلِيَّةِ؟}\\
Are we in an Islamic time or the time of the pre-Islamic ignorance? /hal naħnu fiː zamanin ʔislaːmijjin ʔam zamani ldʒaːhillijjati/

\item \textarabic{هَلْ أَنْتَ فِي عَصْرٍ حَدِيثٍ أَمْ عَصْرِ الْقُرُونِ الْوُسْطَى؟}\\
Are you in a modern era or the era of the Middle Ages? /hal ʔanta fiː ʕasˤrin ħadiːθin ʔam ʕasˤri lquruːni lwustˤaː/

\item \textarabic{هَلْ هَذَا وَقْتٌ فَاطِمِيٌّ أَمْ وَقْتُ الدَّوْلَةِ الْعُثْمَانِيَّةِ؟}\\
Is this a Fatimid time or the time of the Ottoman state? /hal haːðaː waqtun faːtˤimijjun ʔam waqtu ddawlati lʕuθmaːnijjati/

\item \textarabic{هَلْ نَعِيشُ فِي حِقْبَةٍ أَيُّوبِيَّةٍ أَمْ حِقْبَةِ الْمَمَالِيكِ؟}\\
Are we living in an Ayyubid period or the period of the Mamluks? /hal naʕiːʃu fiː ħiqbatin ʔajjuːbijjatin ʔam ħiqbati lmamaːliːki/

\item \textarabic{هَلْ هَذَا عَهْدٌ رَاشِدِيٌّ أَمْ عَهْدُ الْخُلَفَاءِ الْعَبَّاسِيِّينَ؟}\\
Is this a Rashidun epoch or the epoch of the Abbasid caliphs? /hal haːðaː ʕahdun raːʃidijjun ʔam ʕahdu lxulafaːʔi lʕabbaːsijjiːna/

\item \textarabic{هَلْ نَحْنُ بِدَوْلَةٍ مَمْلُوكِيَّةٍ أَمْ دَوْلَةِ الْأَنْدَلُسِ؟}\\
Are we in a Mamluk state or the state of Andalusia? /hal naħnu bidawlatin mamluːkijjatin ʔam dawlati l-ʔandalusi/
\end{enumerate}

% ======================== Levantine Dialect Version ========================
\section{Levantine (Shami) Arabic Dialect}

\begin{tcolorbox}[colback=white,colframe=dialectcolor,title=\textbf{Levantine Version},breakable]
\textarabic{نِحْنَا بِعَصْر عَبَّاسِي وَلَّا عَصْر الدَّوْلِة الْأُمَوِيِّة؟}\\[0.5em]
\textbf{Phonetic:} /niħna biʕasˤar ʕabbaːsiː walla ʕasˤar eddawle l-ʔumawijje/\\[0.5em]
\textbf{Translation:} Are we in an Abbasid era or the era of the Umayyad state?
\end{tcolorbox}

\textbf{Key Dialectal Changes:}
\begin{itemize}
\item \textarabic{هَلْ} → (omitted) — Yes/no questions often indicated by intonation alone in dialect
\item \textarabic{نَحْنُ} → \textarabic{نِحْنَا} /niħna/ — Dialectal pronunciation with final vowel
\item \textarabic{بِعَصْرٍ} → \textarabic{بِعَصْر} /biʕasˤar/ — Loss of case endings (tanwīn)
\item \textarabic{عَبَّاسِيٍّ} → \textarabic{عَبَّاسِي} /ʕabbaːsiː/ — Loss of case endings
\item \textarabic{أَمْ} → \textarabic{وَلَّا} /walla/ — Dialectal form meaning "or" in questions (from \textarabic{وَإِلَّا})
\item \textarabic{عَصْرِ} → \textarabic{عَصْر} /ʕasˤar/ — Loss of case marking
\item \textarabic{الدَّوْلَةِ} → \textarabic{الدَّوْلِة} /eddawle/ — Dialectal pronunciation
\item \textarabic{الْأُمَوِيَّةِ} → \textarabic{الْأُمَوِيِّة} /l-ʔumawijje/ — Simplified pronunciation
\item Overall: Loss of iʕrāb (case system) and some grammatical restructuring
\end{itemize}

\subsection{Alternative Dialectal Versions}

\textbf{Egyptian Arabic:}\\
\textarabic{إِحْنَا فِي عَصْر عَبَّاسِي وَلَّا عَصْر الدَّوْلَة الْأُمَوِيَّة؟}\\
/ʔeħna fi ʕasˤr ʕabbaːsi walla ʕasˤr eddawla l-ʔumawejja/

\textbf{Gulf Arabic:}\\
\textarabic{نِحْنَا بِعَصْر عَبَّاسِي وِلَّا عَصْر الدَّوْلَة الْأُمَوِيَّة؟}\\
/niħna biʕasˤir ʕabbaːsi willa ʕasˤir eddawla l-ʔumawejja/

\textbf{Moroccan Arabic:}\\
\textarabic{حْنَا فْ عَصْر عَبَّاسِي وَلَّا عَصْر الدَّوْلَة الْأُمَوِيَّة؟}\\
/ħna f ʕasˤr ʕabbaːsi wəlla ʕasˤr eddəwla l-ʔumawejja/ 


% ======================== Additional Learning Notes ========================
\section{Additional Learning Notes}
\noindent
\begin{tcolorbox}[breakable,colback=boxcolor,colframe=accentcolor,title=\textbf{Cultural and Historical Context}]

\textbf{Historical Significance:}
\begin{itemize}
\item \textbf{Umayyad Caliphate (661-750 CE):} First hereditary Islamic dynasty, capital Damascus, Arab aristocratic rule, rapid territorial expansion
\item \textbf{Abbasid Caliphate (750-1258 CE):} Overthrew Umayyads, capital Baghdad, multicultural empire, Islamic Golden Age of science and culture
\item The transition from Umayyad to Abbasid marked a shift from Arab tribal dominance to a more cosmopolitan, Persian-influenced administration
\end{itemize}

\textbf{Modern Political Usage:}
This phrase is commonly used rhetorically in modern Arabic political discourse to:
\begin{itemize}
\item Criticize authoritarian or dynastic rule resembling medieval caliphates
\item Question outdated governance models
\item Express frustration with lack of modernization
\item Suggest that current political systems are anachronistic
\item Imply that leaders act like medieval monarchs rather than modern statesmen
\end{itemize}

\textbf{Rhetorical Function:}
The question is typically rhetorical and sarcastic, implying:
\begin{enumerate}
\item Current governance resembles medieval dynasties
\item Political systems haven't evolved since 7th-8th centuries
\item Leaders behave like caliphs rather than elected officials
\item Democratic institutions are absent or ineffective
\end{enumerate}

\textbf{Why These Two Dynasties?}
The Umayyad-Abbasid contrast is meaningful because:
\begin{itemize}
\item Both represent different models of Islamic governance
\item The Abbasid revolution (750 CE) was a major historical rupture
\item They symbolize two distinct political philosophies
\item Both ended centuries ago, emphasizing anachronism
\item They're well-known reference points in Arab historical consciousness
\end{itemize}

\textbf{Grammatical Note on Asymmetry:}
The phrase shows interesting asymmetry:
\begin{itemize}
\item First option: \textarabic{بِعَصْرٍ عَبَّاسِيٍّ} (indefinite) — "an Abbasid era" (generic)
\item Second option: \textarabic{عَصْرِ الدَّوْلَةِ الْأُمَوِيَّةِ} (definite) — "the era of the Umayyad state" (specific)
\item This asymmetry may reflect historical perspective: multiple Abbasid periods vs. the singular Umayyad caliphate
\end{itemize}
\end{tcolorbox}

\subsection{Memory Tips}

\begin{enumerate}
\item \textbf{Question Pattern:} \textarabic{هَلْ + pronoun + predicates with أَمْ} is standard for yes/no alternative questions

\item \textbf{Dynasty Names:} Remember nisba formation:
\begin{itemize}
\item \textarabic{عَبَّاس} (ʕAbbās) + \textarabic{ـِيّ} → \textarabic{عَبَّاسِيّ} (Abbasid)
\item \textarabic{أُمَيَّة} (Umayya) + \textarabic{ـِيّ} → \textarabic{أُمَوِيّ} (Umayyad, note the \textarabic{و} insertion)
\end{itemize}

\item \textbf{Chronological Order:} Umayyad (661-750) came BEFORE Abbasid (750-1258)
\begin{itemize}
\item Mnemonic: "U before A" (Umayyad before Abbasid)
\end{itemize}

\item \textbf{Capital Cities:} 
\begin{itemize}
\item Umayyad = Damascus (both start with "D" sound in Arabic: \textarabic{أُمَوِيّ/دِمَشْق})
\item Abbasid = Baghdad (both have "b" sound: \textarabic{عَبَّاسِيّ/بَغْدَاد})
\end{itemize}

\item \textbf{Preposition \textarabic{بِ}:} When expressing "in an era," \textarabic{بِ} is more common than \textarabic{فِي} with \textarabic{عَصْر}

\item \textbf{Alternative Questions:} Remember: \textarabic{أَمْ} for questions, \textarabic{أَوْ} for statements

\item \textbf{Ellipsis Recognition:} Second predicate omits \textarabic{بِ} — this is normal in coordinated structures
\end{enumerate}

\subsection{Related Vocabulary Families}

\begin{multicols}{2}

\textbf{Time/Era Words:}\\
\textarabic{عَصْر} — era /ʕasˤr/\\
\textarabic{زَمَن} — time /zaman/\\
\textarabic{عَهْد} — epoch /ʕahd/\\
\textarabic{دَهْر} — age /dahr/\\
\textarabic{حِقْبَة} — period /ħiqba/\\
\textarabic{وَقْت} — time /waqt/\\
\textarabic{أَوَان} — time /ʔawaːn/\\
\textarabic{فَتْرَة} — period /fatra/\\

\textbf{State/Dynasty Words:}\\
\textarabic{دَوْلَة} — state /dawla/\\
\textarabic{خِلَافَة} — caliphate /xilaːfa/\\
\textarabic{إِمْبَرَاطُورِيَّة} — empire /ʔimbraːtˤuːrijja/\\
\textarabic{مَمْلَكَة} — kingdom /mamlaka/\\
\textarabic{سُلْطَنَة} — sultanate /sultˤana/\\
\textarabic{حُكُومَة} — government /ħukuːma/\\
\textarabic{نِظَام} — regime /nidˤaːm/\\
\textarabic{إِمَارَة} — emirate /ʔimaːra/\\

\textbf{Islamic Dynasties:}\\
\textarabic{أُمَوِيّ} — Umayyad /ʔumawijj/\\
\textarabic{عَبَّاسِيّ} — Abbasid /ʕabbaːsijj/\\
\textarabic{فَاطِمِيّ} — Fatimid /faːtˤimijj/\\
\textarabic{أَيُّوبِيّ} — Ayyubid /ʔajjuːbijj/\\
\textarabic{مَمْلُوكِيّ} — Mamluk /mamluːkijj/\\
\textarabic{عُثْمَانِيّ} — Ottoman /ʕuθmaːnijj/\\
\textarabic{رَاشِدِيّ} — Rashidun /raːʃidijj/\\
\textarabic{سَلْجُوقِيّ} — Seljuk /saldʒuːqijj/\\

\textbf{Question Words:}\\
\textarabic{هَلْ} — (yes/no) /hal/\\
\textarabic{مَا} — what /maː/\\
\textarabic{مَنْ} — who /man/\\
\textarabic{مَتَى} — when /mataː/\\
\textarabic{أَيْنَ} — where /ʔajna/\\
\textarabic{كَيْفَ} — how /kajfa/\\
\textarabic{لِمَاذَا} — why /limaːðaː/\\
\textarabic{كَمْ} — how many /kam/\\

\textbf{Conjunctions:}\\
\textarabic{أَمْ} — or (questions) /ʔam/\\
\textarabic{أَوْ} — or (statements) /ʔaw/\\
\textarabic{وَ} — and /wa/\\
\textarabic{فَ} — so, then /fa/\\
\textarabic{لَكِنْ} — but /laːkin/\\
\textarabic{بَلْ} — rather /bal/\\
\textarabic{إِمَّا} — either /ʔimmaː/\\
\textarabic{ثُمَّ} — then /θumma/\\

\textbf{Historical Periods:}\\
\textarabic{الْجَاهِلِيَّة} — pre-Islamic era /aldʒaːhillijja/\\
\textarabic{صَدْر الْإِسْلَام} — early Islam /sˤadr al-ʔislaːm/\\
\textarabic{الْعَصْر الذَّهَبِيّ} — Golden Age /alʕasˤr aðːðahabijj/\\
\textarabic{الْقُرُون الْوُسْطَى} — Middle Ages /alquruːn alwustˤaː/\\
\textarabic{عَصْر النَّهْضَة} — Renaissance /ʕasˤr annahðˤa/\\
\textarabic{الْعَصْر الْحَدِيث} — modern era /alʕasˤr alħadiːθ/\\
\end{multicols}

\section{Grammatical Deep Dive}

\subsection{Yes/No Questions with \textarabic{هَلْ}}

\begin{tcolorbox}[colback=boxcolor,colframe=headercolor,breakable]
\textbf{Function:} \textarabic{هَلْ} introduces polar questions requiring yes/no answers

\textbf{Structure Patterns:}

\begin{tabular}{|l|p{10cm}|}
\hline
\textbf{Pattern} & \textbf{Example} \\
\hline
\textarabic{هَلْ} + nominal sentence & \textarabic{هَلْ أَنْتَ طَالِبٌ؟} — Are you a student? \\
\hline
\textarabic{هَلْ} + verbal sentence & \textarabic{هَلْ ذَهَبْتَ إِلَى الْمَدْرَسَةِ؟} — Did you go to school? \\
\hline
\textarabic{هَلْ} + alternatives & \textarabic{هَلْ هَذَا أَمْ ذَاكَ؟} — Is it this or that? \\
\hline
\end{tabular}

\vspace{0.5em}
\textbf{Expected Responses:}
\begin{itemize}
\item Affirmative: \textarabic{نَعَمْ} (yes), \textarabic{أَجَلْ} (yes, indeed), \textarabic{بَلَى} (yes, contradicting negative)
\item Negative: \textarabic{لَا} (no), \textarabic{كَلَّا} (no, certainly not)
\end{itemize}

\vspace{0.5em}
\textbf{Comparison with Other Interrogatives:}
\begin{itemize}
\item \textarabic{هَلْ} — yes/no questions
\item \textarabic{أَ} — general interrogative particle (less common in MSA)
\item \textarabic{مَا} — what questions
\item \textarabic{مَنْ} — who questions
\end{itemize}
\end{tcolorbox}

\subsection{Alternative Questions with \textarabic{أَمْ}}

\begin{tcolorbox}[colback=white,colframe=accentcolor,breakable]
\textbf{Two Types of \textarabic{أَمْ} Questions:}


\textbf{1. Connected Alternative (\textarabic{أَمْ الْمُتَّصِلَة}):}\\
Used with initial \textarabic{هَلْ} or \textarabic{أَ}, offers two alternatives:
\begin{itemize}
\item \textarabic{هَلْ نَحْنُ بِعَصْرٍ عَبَّاسِيٍّ أَمْ عَصْرِ الدَّوْلَةِ الْأُمَوِيَّةِ؟}
\item Are we in an Abbasid era or (in) the era of the Umayyad state?
\end{itemize}

\textbf{2. Disconnected Alternative (\textarabic{أَمْ الْمُنْقَطِعَة}):}\\
Used independently, means "rather, but rather":
\begin{itemize}
\item \textarabic{هَلْ نَجَحْتَ، أَمْ أَنْتَ لَمْ تَدْرُسْ؟}
\item Did you succeed, or did you not study?
\end{itemize}

\vspace{0.5em}
\textbf{Ellipsis with \textarabic{أَمْ}:} \\
Common to omit repeated elements in the second alternative:
\begin{itemize}
\item Full: \textarabic{هَلْ نَحْنُ بِعَصْرٍ عَبَّاسِيٍّ أَمْ [نَحْنُ] بِعَصْرِ الدَّوْلَةِ الْأُمَوِيَّةِ؟}
\item Actual: \textarabic{هَلْ نَحْنُ بِعَصْرٍ عَبَّاسِيٍّ أَمْ عَصْرِ الدَّوْلَةِ الْأُمَوِيَّةِ؟}
\end{itemize}
\end{tcolorbox}

\subsection{Nisba Adjectives: Formation and Usage}

\begin{tcolorbox}[colback=boxcolor,colframe=headercolor,breakable]
\textbf{Definition:} Nisba adjectives (\textarabic{النِّسْبَة}) relate nouns to places, tribes, times, or abstract concepts by adding suffix \textarabic{ـِيّ} /-ijj/ (masc.) or \textarabic{ـِيَّة} /-ijja/ (fem.)

\textbf{Formation Rules:}

\begin{tabular}{|l|l|l|}
\hline
\textbf{Base Type} & \textbf{Process} & \textbf{Example} \\
\hline
Regular noun & Add \textarabic{ـِيّ} & \textarabic{مِصْر} → \textarabic{مِصْرِيّ} (Egyptian) \\
\hline
Tāʔ marbūṭah & Remove \textarabic{ة}, add \textarabic{ـِيّ} & \textarabic{مَكَّة} → \textarabic{مَكِّيّ} (Meccan) \\
\hline
Alif maqṣūrah & Change \textarabic{ى} to \textarabic{وِيّ} & \textarabic{فَتْوَى} → \textarabic{فَتْوَوِيّ} (advisory) \\
\hline
Hollow verb & Insert \textarabic{و} or \textarabic{ي} & \textarabic{أُمَيَّة} → \textarabic{أُمَوِيّ} (Umayyad) \\
\hline
\end{tabular}

\vspace{0.5em}
\textbf{Historical Dynasty Nisbas:}

\begin{tabular}{|l|l|l|}
\hline
\textbf{Founder/Ancestor} & \textbf{Nisba} & \textbf{English} \\
\hline
\textarabic{عَبَّاس} & \textarabic{عَبَّاسِيّ} & Abbasid \\
\hline
\textarabic{أُمَيَّة} & \textarabic{أُمَوِيّ} & Umayyad \\
\hline
\textarabic{فَاطِمَة} & \textarabic{فَاطِمِيّ} & Fatimid \\
\hline
\textarabic{أَيُّوب} & \textarabic{أَيُّوبِيّ} & Ayyubid \\
\hline
\textarabic{عُثْمَان} & \textarabic{عُثْمَانِيّ} & Ottoman \\
\hline
\textarabic{رَشِيد} & \textarabic{رَاشِدِيّ} & Rashidun \\
\hline
\end{tabular}

\vspace{0.5em}
\textbf{Agreement Pattern:}\\
Nisba adjectives follow standard adjective agreement rules:

\begin{tabular}{|l|l|l|}
\hline
\textbf{Modified Noun} & \textbf{Nisba Form} & \textbf{Example} \\
\hline
\textarabic{عَصْرٌ} (masc. sing.) & \textarabic{عَبَّاسِيٌّ} & \textarabic{عَصْرٌ عَبَّاسِيٌّ} \\
\hline
\textarabic{دَوْلَةٌ} (fem. sing.) & \textarabic{عَبَّاسِيَّةٌ} & \textarabic{دَوْلَةٌ عَبَّاسِيَّةٌ} \\
\hline
\textarabic{عُصُورٌ} (plural) & \textarabic{عَبَّاسِيَّةٌ} & \textarabic{عُصُورٌ عَبَّاسِيَّةٌ} \\
\hline
\end{tabular}
\end{tcolorbox}

\subsection{Iḍāfah Construction: Definite vs. Indefinite}

\begin{tcolorbox}[colback=white,colframe=accentcolor,breakable]
\textbf{Comparison of Two Structures in Our Phrase:}

\textbf{Structure 1: Indefinite}\\
\textarabic{بِعَصْرٍ عَبَّاسِيٍّ} — "in an Abbasid era"

\begin{itemize}
\item \textbf{Type:} Prepositional phrase + indefinite noun + indefinite adjective
\item \textbf{Meaning:} Generic reference to any Abbasid period
\item \textbf{Case marking:} Both noun and adjective show tanwīn
\item \textbf{Definiteness:} Completely indefinite
\end{itemize}

\textbf{Structure 2: Definite Iḍāfah}\\
\textarabic{عَصْرِ الدَّوْلَةِ الْأُمَوِيَّةِ} — "the era of the Umayyad state"

\begin{itemize}
\item \textbf{Type:} Iḍāfah construction with definite second term
\item \textbf{Components:}
  \begin{itemize}
  \item Muḍāf: \textarabic{عَصْرِ} (first term, no \textarabic{ال}, no tanwīn)
  \item Muḍāf ilayh: \textarabic{الدَّوْلَةِ} (second term, definite)
  \item Adjective: \textarabic{الْأُمَوِيَّةِ} (agrees with \textarabic{الدَّوْلَةِ})
  \end{itemize}
\item \textbf{Meaning:} Specific reference to the historical Umayyad state
\item \textbf{Definiteness:} Whole construction is definite (through second term)
\end{itemize}

\vspace{0.5em}
\textbf{Key Iḍāfah Rules Illustrated:}
\begin{enumerate}
\item First term (\textarabic{عَصْرِ}) never takes \textarabic{ال} or tanwīn
\item First term's case depends on syntactic role (here: genitive after \textarabic{بِ})
\item Second term (\textarabic{الدَّوْلَةِ}) must be genitive
\item Second term's definiteness determines whole construction's definiteness
\item Adjectives agree with the noun they modify in all features
\end{enumerate}

\vspace{0.5em}
\textbf{Adjective Placement Options:}

\begin{tabular}{|l|p{6cm}|l|}
\hline
\textbf{Position} & \textbf{Example} & \textbf{What it modifies} \\
\hline
After entire iḍāfah & \textarabic{عَصْرُ الدَّوْلَةِ الْكَبِيرُ} & First term (\textarabic{عَصْر}) \\
\hline
After second term & \textarabic{عَصْرُ الدَّوْلَةِ الْكَبِيرَةِ} & Second term (\textarabic{دَوْلَة}) \\
\hline
\end{tabular}

In our phrase: \textarabic{الْأُمَوِيَّةِ} modifies \textarabic{الدَّوْلَةِ} (feminine, genitive, definite)
\end{tcolorbox}

\subsection{Prepositional Phrases in Predicates}

\begin{tcolorbox}[colback=boxcolor,colframe=headercolor,breakable]
\textbf{Using Prepositions as Predicates:}

In nominal sentences, prepositional phrases commonly function as predicates:

\begin{tabular}{|l|p{5cm}|p{5cm}|}
\hline
\textbf{Structure} & \textbf{Arabic} & \textbf{English} \\
\hline
Subject + \textarabic{فِي} + noun & \textarabic{الْكِتَابُ فِي الْحَقِيبَةِ} & The book is in the bag \\
\hline
Subject + \textarabic{عَلَى} + noun & \textarabic{الْقَلَمُ عَلَى الطَّاوِلَةِ} & The pen is on the table \\
\hline
Subject + \textarabic{بِ} + noun & \textarabic{نَحْنُ بِعَصْرٍ جَدِيدٍ} & We are in a new era \\
\hline
\end{tabular}

\vspace{0.5em}
\textbf{Why \textarabic{بِ} Instead of \textarabic{فِي}?}

Both can mean "in," but with nuances:
\begin{itemize}
\item \textarabic{فِي} — more literal location: \textarabic{فِي الْبَيْتِ} (in the house)
\item \textarabic{بِ} — more abstract, metaphorical: \textarabic{بِعَصْرٍ} (in an era/circumstances)
\item \textarabic{بِ} — emphasis on circumstance or state: \textarabic{بِحَالَةٍ جَيِّدَةٍ} (in good condition)
\end{itemize}

\vspace{0.5em}
\textbf{Common Prepositions Governing Genitive:}

\begin{multicols}{3}
\textarabic{بِ} /bi/ — with, in\\
\textarabic{مِنْ} /min/ — from\\
\textarabic{إِلَى} /ʔilaː/ — to\\
\textarabic{عَلَى} /ʕalaː/ — on\\
\textarabic{فِي} /fiː/ — in\\
\textarabic{عَنْ} /ʕan/ — about\\
\textarabic{لِ} /li/ — for\\
\textarabic{كَ} /ka/ — like\\
\textarabic{مَعَ} /maʕa/ — with\\
\end{multicols}

All prepositions require the following noun to be in genitive case.
\end{tcolorbox}

\section{Practice Exercises}

\subsection{Translation Practice}

Translate the following into Arabic using similar structures:

\begin{enumerate}
\item Are we in a modern era or the era of the Middle Ages?
\vspace{2em}

\item Are you in a Fatimid state or the state of the Mamluks?
\vspace{2em}

\item Is this an Ottoman period or the period of the Rashidun Caliphate?
\vspace{2em}

\item Are they in an Islamic time or the time of ignorance?
\vspace{2em}
\end{enumerate}

\subsection{Grammatical Transformation}

Transform the original phrase in the following ways:

\begin{enumerate}
\item Change to past tense using \textarabic{كَانَ}:
\vspace{2em}

\item Change subject from "we" to "you (masc. pl.)":
\vspace{2em}

\item Make both alternatives indefinite:
\vspace{2em}

\item Add temporal context (yesterday, today, tomorrow):
\vspace{2em}
\end{enumerate}

\subsection{Nisba Formation Practice}

Form nisba adjectives from the following nouns:

\begin{enumerate}
\item \textarabic{مِصْر} (Egypt) → \underline{\hspace{3cm}}

\item \textarabic{سُورِيَا} (Syria) → \underline{\hspace{3cm}}

\item \textarabic{الْعِرَاق} (Iraq) → \underline{\hspace{3cm}}

\item \textarabic{فَاطِمَة} (Fatima) → \underline{\hspace{3cm}}

\item \textarabic{الْأَنْدَلُس} (Andalusia) → \underline{\hspace{3cm}}
\end{enumerate}

\subsection{Case Analysis}

Identify the grammatical case of each underlined word and explain why:

\begin{enumerate}
\item \textarabic{هَلْ نَحْنُ بِ\underline{عَصْرٍ} \underline{عَبَّاسِيٍّ}؟}\\
Case and reason: \underline{\hspace{8cm}}

\item \textarabic{عَصْرِ \underline{الدَّوْلَةِ} \underline{الْأُمَوِيَّةِ}}\\
Case and reason: \underline{\hspace{8cm}}

\item \textarabic{هَلْ \underline{نَحْنُ} فِي بَيْتٍ؟}\\
Case and reason: \underline{\hspace{8cm}}
\end{enumerate}

\section{Answers to Practice Exercises}

\subsection{Translation Practice — Suggested Answers}

\begin{enumerate}
\item \textarabic{هَلْ نَحْنُ بِعَصْرٍ حَدِيثٍ أَمْ عَصْرِ الْقُرُونِ الْوُسْطَى؟}\\
/hal naħnu biʕasˤrin ħadiːθin ʔam ʕasˤri lquruːni lwustˤaː/

\item \textarabic{هَلْ أَنْتَ بِدَوْلَةٍ فَاطِمِيَّةٍ أَمْ دَوْلَةِ الْمَمَالِيكِ؟}\\
/hal ʔanta bidawlatin faːtˤimijjatin ʔam dawlati lmamaːliːki/

\item \textarabic{هَلْ هَذَا عَهْدٌ عُثْمَانِيٌّ أَمْ عَهْدُ الْخِلَافَةِ الرَّاشِدَةِ؟}\\
/hal haːðaː ʕahdun ʕuθmaːnijjun ʔam ʕahdu lxilaːfati rraːʃidati/

\item \textarabic{هَلْ هُمْ فِي زَمَنٍ إِسْلَامِيٍّ أَمْ زَمَنِ الْجَاهِلِيَّةِ؟}\\
/hal hum fiː zamanin ʔislaːmijjin ʔam zamani ldʒaːhillijjati/
\end{enumerate}

\subsection{Grammatical Transformation — Answers}

\begin{enumerate}
\item \textbf{Past tense with \textarabic{كَانَ}:}\\
\textarabic{هَلْ كُنَّا بِعَصْرٍ عَبَّاسِيٍّ أَمْ عَصْرِ الدَّوْلَةِ الْأُمَوِيَّةِ؟}\\
/hal kunnaː biʕasˤrin ʕabbaːsijjin ʔam ʕasˤri ddawlati l-ʔumawijjati/\\
Were we in an Abbasid era or the era of the Umayyad state?

\item \textbf{Subject "you" (masc. pl.):}\\
\textarabic{هَلْ أَنْتُمْ بِعَصْرٍ عَبَّاسِيٍّ أَمْ عَصْرِ الدَّوْلَةِ الْأُمَوِيَّةِ؟}\\
/hal ʔantum biʕasˤrin ʕabbaːsijjin ʔam ʕasˤri ddawlati l-ʔumawijjati/\\
Are you (pl.) in an Abbasid era or the era of the Umayyad state?

\item \textbf{Both alternatives indefinite:}\\
\textarabic{هَلْ نَحْنُ بِعَصْرٍ عَبَّاسِيٍّ أَمْ عَصْرٍ أُمَوِيٍّ؟}\\
/hal naħnu biʕasˤrin ʕabbaːsijjin ʔam ʕasˤrin ʔumawijjin/\\
Are we in an Abbasid era or an Umayyad era?

\item \textbf{With temporal context:}\\
\textarabic{هَلْ نَحْنُ الْيَوْمَ بِعَصْرٍ عَبَّاسِيٍّ أَمْ عَصْرِ الدَّوْلَةِ الْأُمَوِيَّةِ؟}\\
/hal naħnu ljawma biʕasˤrin ʕabbaːsijjin ʔam ʕasˤri ddawlati l-ʔumawijjati/\\
Are we today in an Abbasid era or the era of the Umayyad state?
\end{enumerate}

\subsection{Nisba Formation — Answers}

\begin{enumerate}
\item \textarabic{مِصْر} → \textbf{\textarabic{مِصْرِيّ}} /misˤrijj/ (Egyptian)

\item \textarabic{سُورِيَا} → \textbf{\textarabic{سُورِيّ}} /suːrijj/ (Syrian)

\item \textarabic{الْعِرَاق} → \textbf{\textarabic{عِرَاقِيّ}} /ʕiraːqijj/ (Iraqi)

\item \textarabic{فَاطِمَة} → \textbf{\textarabic{فَاطِمِيّ}} /faːtˤimijj/ (Fatimid)

\item \textarabic{الْأَنْدَلُس} → \textbf{\textarabic{أَنْدَلُسِيّ}} /ʔandalusijj/ (Andalusian)
\end{enumerate}

\subsection{Case Analysis — Answers}

\begin{enumerate}
\item \textarabic{عَصْرٍ} and \textarabic{عَبَّاسِيٍّ}: \textbf{Genitive (majrūr)}\\
Reason: Both are in genitive case because they follow the preposition \textarabic{بِ}. The adjective \textarabic{عَبَّاسِيٍّ} agrees with \textarabic{عَصْرٍ} in case, gender, number, and definiteness.

\item \textarabic{الدَّوْلَةِ} and \textarabic{الْأُمَوِيَّةِ}: \textbf{Genitive (majrūr)}\\
Reason: \textarabic{الدَّوْلَةِ} is genitive as the second term (muḍāf ilayh) of the iḍāfah construction. \textarabic{الْأُمَوِيَّةِ} is genitive because it agrees with \textarabic{الدَّوْلَةِ}.

\item \textarabic{نَحْنُ}: \textbf{Nominative (marfūʕ)}\\
Reason: As an independent pronoun serving as the subject (mubtadaʔ) of a nominal sentence, it is in the nominative case (though pronouns don't show case inflection, they are considered nominative by position).
\end{enumerate}

\section{Extended Historical Context}

\begin{tcolorbox}[colback=boxcolor,colframe=accentcolor,title=\textbf{Detailed Historical Background},breakable]

\subsection*{The Umayyad Caliphate (661-750 CE)}

\textbf{Key Features:}
\begin{itemize}
\item \textbf{Founder:} Muʕāwiyah ibn Abī Sufyān
\item \textbf{Capital:} Damascus (\textarabic{دِمَشْق})
\item \textbf{Territory:} Expanded from Spain to Central Asia
\item \textbf{Administration:} Arab aristocratic rule, hereditary succession
\item \textbf{Language:} Arabization of administration
\item \textbf{Legacy:} Spread of Islam and Arabic language
\end{itemize}

\textbf{Notable Rulers:}
\begin{itemize}
\item Muʕāwiyah I (661-680) — Founder
\item ʕAbd al-Malik (685-705) — Arabization, Dome of the Rock
\item al-Walīd I (705-715) — Greatest expansion, Umayyad Mosque
\item ʕUmar II (717-720) — Pious ruler, reforms
\item Hishām (724-743) — Administrative consolidation
\end{itemize}

\subsection*{The Abbasid Caliphate (750-1258 CE)}

\textbf{Key Features:}
\begin{itemize}
\item \textbf{Founder:} Abū al-ʕAbbās as-Saffāḥ
\item \textbf{Capital:} Baghdad (\textarabic{بَغْدَاد}) — "City of Peace"
\item \textbf{Territory:} Initially similar to Umayyad, gradually fragmented
\item \textbf{Administration:} Persian-influenced, bureaucratic, multicultural
\item \textbf{Culture:} Islamic Golden Age — science, philosophy, literature
\item \textbf{Legacy:} Cultural and intellectual flourishing
\end{itemize}

\textbf{Notable Rulers:}
\begin{itemize}
\item Abū Jaʕfar al-Manṣūr (754-775) — Founded Baghdad
\item Hārūn ar-Rashīd (786-809) — Golden Age apex, "Arabian Nights"
\item al-Maʔmūn (813-833) — House of Wisdom, translation movement
\item al-Muʕtaṣim (833-842) — Military reforms, Turkish guards
\item al-Mutawakkil (847-861) — Last strong caliph before fragmentation
\end{itemize}

\subsection*{The Abbasid Revolution (750 CE)}

\textbf{Causes:}
\begin{itemize}
\item Resentment of Umayyad Arab supremacy
\item Persian (mawālī) dissatisfaction with second-class status
\item Religious opposition from Shīʕa and Khārijites
\item Economic grievances in outlying provinces
\item Propaganda from Abbasid movement in Khurasan
\end{itemize}

\textbf{Consequences:}
\begin{itemize}
\item Near-complete elimination of Umayyad family
\item Shift of power from Syria to Iraq
\item Increased Persian cultural influence
\item More inclusive, less Arab-centric empire
\item Foundation for Islamic Golden Age
\end{itemize}

\subsection*{Modern Political Usage}

This phrase has become a rhetorical device in contemporary Arab political discourse:

\textbf{Common contexts:}
\begin{itemize}
\item Criticizing hereditary rule and dynastic politics
\item Questioning lack of democratic institutions
\item Highlighting outdated governance models
\item Expressing frustration with authoritarian leadership
\item Comparing modern states to medieval caliphates
\end{itemize}

\textbf{Implied criticism:}
\begin{itemize}
\item Political systems haven't evolved in 1,300+ years
\item Leaders behave like absolute monarchs
\item Governance resembles tribal/dynastic models
\item Absence of modern democratic principles
\item Succession based on family, not merit or election
\end{itemize}
\end{tcolorbox}

\section{Additional Linguistic Notes}

\subsection{Phonological Features}

\begin{tcolorbox}[colback=white,colframe=accentcolor,breakable]
\textbf{Important Pronunciation Points:}

\begin{enumerate}
\item \textbf{Pharyngealized consonants:}
\begin{itemize}
\item \textarabic{ع} /ʕ/ — pharyngeal fricative (not in English)
\item \textarabic{ص} /sˤ/ — emphatic s (in \textarabic{عَصْر})
\end{itemize}

\item \textbf{Doubled consonants in nisba:}
\begin{itemize}
\item \textarabic{عَبَّاسِيّ} — note gemination of /b/ and /j/
\item Must pronounce both instances distinctly
\end{itemize}

\item \textbf{Case ending pronunciation:}
\begin{itemize}
\item Formal MSA: \textarabic{عَصْرٍ عَبَّاسِيٍّ} /ʕasˤrin ʕabbaːsijjin/
\item Pausal form: \textarabic{عَصْر عَبَّاسِيّ} /ʕasˤr ʕabbaːsijj/
\item Modern reading often omits short final vowels
\end{itemize}

\item \textbf{Definite article assimilation:}
\begin{itemize}
\item \textarabic{الدَّوْلَة} /addawla/ not */aldawla/
\item \textarabic{د} is a "sun letter" — causes assimilation
\item \textarabic{الْأُمَوِيَّة} /al-ʔumawijja/ — \textarabic{أ} is "moon letter"
\end{itemize}
\end{enumerate}
\end{tcolorbox}

\subsection{Stylistic Variations}

\begin{tcolorbox}[colback=boxcolor,colframe=headercolor,breakable]
\textbf{Alternative Ways to Express the Same Question:}

\begin{enumerate}
\item \textbf{More formal/literary:}\\
\textarabic{أَنَحْنُ فِي عَصْرٍ عَبَّاسِيٍّ أَمْ فِي عَصْرِ الدَّوْلَةِ الْأُمَوِيَّةِ؟}\\
(Using \textarabic{أَ} instead of \textarabic{هَلْ}, repeating \textarabic{فِي})

\item \textbf{With emphasis:}\\
\textarabic{هَلْ نَحْنُ حَقًّا بِعَصْرٍ عَبَّاسِيٍّ أَمْ عَصْرِ الدَّوْلَةِ الْأُمَوِيَّةِ؟}\\
(Adding \textarabic{حَقًّا} "really, truly")

\item \textbf{Rhetorical intensification:}\\
\textarabic{فِي أَيِّ عَصْرٍ نَحْنُ؟ عَصْرٍ عَبَّاسِيٍّ أَمْ عَصْرِ الدَّوْلَةِ الْأُمَوِيَّةِ؟}\\
(In which era are we? An Abbasid era or the era of the Umayyad state?)

\item \textbf{With temporal marker:}\\
\textarabic{هَلْ مَا زِلْنَا فِي عَصْرٍ عَبَّاسِيٍّ أَمْ عَصْرِ الدَّوْلَةِ الْأُمَوِيَّةِ؟}\\
(Are we still in an Abbasid era or the era of the Umayyad state?)

\item \textbf{Negative form:}\\
\textarabic{أَلَسْنَا بِعَصْرٍ حَدِيثٍ بَلْ بِعَصْرٍ عَبَّاسِيٍّ أَوْ أُمَوِيٍّ؟}\\
(Are we not in a modern era but rather in an Abbasid or Umayyad era?)
\end{enumerate}
\end{tcolorbox}

%\clearpage
\section{Final Summary}
\noindent
\begin{tcolorbox}[colback=boxcolor,colframe=accentcolor,breakable]
\textbf{Key Takeaways from This Phrase:}

\textbf{Grammatical Structures Mastered:}
\begin{itemize}
\item Yes/no questions with \textarabic{هَلْ}
\item Alternative questions with \textarabic{أَمْ}
\item Prepositional phrases as predicates
\item Iḍāfah constructions (definite and indefinite)
\item Nisba adjective formation and agreement
\item Ellipsis in coordinate structures
\item Case assignment in complex phrases
\end{itemize}

\textbf{Vocabulary Domains:}
\begin{itemize}
\item Time and era vocabulary
\item Islamic dynasties and historical periods
\item Political and governance terminology
\item Question formation patterns
\item Prepositional usage
\end{itemize}

\textbf{Cultural Knowledge:}
\begin{itemize}
\item Umayyad vs. Abbasid historical context
\item Modern rhetorical usage in political discourse
\item Understanding of anachronism as criticism
\item Arab historical consciousness
\item Contemporary political commentary
\end{itemize}

\textbf{Practical Application:}
\begin{itemize}
\item Understanding modern Arab political rhetoric
\item Forming alternative questions in Arabic
\item Using historical references metaphorically
\item Constructing complex nominal sentences
\item Applying nisba adjectives correctly
\end{itemize}
\end{tcolorbox}

\end{document}